%---------------------------------------------------------------------------
% Dennis Johan Loevlie -- Academic CV (canonical)
% IBM Plex stack, Tier 2-6 polish, hyperxmp + bookmark + xurl,
% Research Interests / Open Source / Engineering Skills sections.
% NB: We use hyperxmp (XMP packet that Adobe Bridge / Google Scholar /
% ATS pipelines actually parse). \DocumentMetadata{pdfstandard=ua-2}
% would conflict with hyperxmp, so we pick the higher-ROI option for
% ATS readability over PDF/UA-2 accessibility.
%---------------------------------------------------------------------------
\documentclass[10pt,letterpaper]{article}

% cmap must precede fontenc to apply Unicode mapping to all loaded fonts.
\usepackage{cmap}
\usepackage[T1]{fontenc}
\usepackage[utf8]{inputenc}
\usepackage[margin=0.55in,top=0.45in,bottom=0.5in]{geometry}

% --- Typography ----------------------------------------------------------
\usepackage[scale=0.97]{plex-serif}
\usepackage[scale=0.91]{plex-sans}
\usepackage[scale=0.89]{plex-mono}
\usepackage[protrusion=true,expansion=true,tracking=true,kerning=true,final]{microtype}

% Old-style proportional figures by default in body text — Bringhurst
% rule: numbers in running prose should sit at x-height, not cap-height,
% so they integrate with lowercase letterforms. Plex Serif's text-osf
% variant is `TOsF`. Single-arg `osf` selector avoids the missing
% prop+osf combo that breaks Plex's pdflatex font lookup.
\AtBeginDocument{\figureversion{osf}}

% Single uniform space after periods (modern norm)
\frenchspacing

% Family helpers
\newcommand{\sansfam}{\sffamily}
\newcommand{\monofam}{\ttfamily}

\linespread{1.06}
\setlength{\parindent}{0pt}

\usepackage{titlesec}
\usepackage{enumitem}
\usepackage[dvipsnames]{xcolor}
\usepackage{parskip}
% Bringhurst vertical rhythm: paragraph spacing as a fraction of the
% baseline so the grid is consistent. 0.25× baseline keeps density tight
% on a 2-page target while snapping to the grid the eye reads as "designed."
\setlength{\parskip}{0.25\baselineskip plus 0.05\baselineskip minus 0.05\baselineskip}
\usepackage{xspace}

% Palette — defined before hyperref so accentlight is in scope for link colors.
% \colorlet tint ramp gives a coherent scale from the same accent base
% rather than picking arbitrary hex codes.
\definecolor{accent}{HTML}{003153}  % Prussian blue
\definecolor{accentlight}{HTML}{1E5288}  % Prussian light
\definecolor{body}{HTML}{1A1A1A}
\definecolor{meta}{HTML}{6B6B6B}
\definecolor{rule}{HTML}{C8C8C8}
\colorlet{accent80}{accent!80!white}
\colorlet{accent50}{accent!50!white}
\colorlet{accent20}{accent!20!white}
\color{body}

% Hyperref with full PDF metadata. hyperxmp adds an XMP packet that
% Adobe Bridge, Google Scholar, JSTOR, and ATS pipelines actually read
% (legacy Info dict alone misses these). bookmark replaces hyperref's
% bookmark machinery — first-pass writes, color/styling support, and
% tolerates skipped levels in the outline tree. xurl enables long-URL
% line breaks so a deep DOI doesn't overflow the column.
\usepackage[
  pdfusetitle,
  pdfauthor={Dennis Johan Loevlie},
  pdftitle={Curriculum Vitae --- Dennis Johan Loevlie},
  pdfsubject={Academic CV},
  pdfkeywords={machine learning, multimodal ML, multimodal foundation models, tabular foundation models, multiple instance learning, RL post-training, computer vision, medical imaging, probabilistic deep learning},
  pdflang=en-US,
  hidelinks,
  bookmarks=true, bookmarksnumbered=false, bookmarksopen=true, bookmarksopenlevel=1,
  colorlinks=true,
  urlcolor=accentlight, linkcolor=accentlight, citecolor=accentlight
]{hyperref}
\usepackage{hyperxmp}
\hypersetup{
  pdfauthortitle={ML Researcher},
  pdfcontactemail={dennis.loevlie@tufts.edu},
  pdfcontacturl={https://dennisloevlie.com},
  pdfcopyright={Copyright (c) 2026 Dennis Loevlie. All rights reserved.}
}
\usepackage{bookmark}
\bookmarksetup{color=accent, bold, open=true, openlevel=1, numbered=false}
\usepackage{xurl}
\input{glyphtounicode}\pdfgentounicode=1
\usepackage{fancyhdr}
\usepackage{lastpage}
\usepackage{needspace}

% --- Bibliography (biblatex + biblatex-publist) --------------------------
% biblatex-publist is purpose-built for academic CV publication lists:
% omits-name capability, year-clusters, native author+an annotations
% for highlighting your name, and clean keyword-filter for "Selected"
% vs full lists. Build pipeline: pdflatex → biber → pdflatex × 2.
\usepackage[
  backend   = biber,
  style     = numeric,
  sorting   = ydnt,
  maxnames  = 99,
  minnames  = 1,
  giveninits= true,
  uniquename= init,
  doi       = false,
  isbn      = false,
  url       = false
]{biblatex}
\addbibresource{publications.bib}

% Bold the highlighted-name parts (per author+an = {N=highlight} in .bib).
% Per Bringhurst rule 3.2.1 we don't double-emphasize — bold only.
\renewcommand*{\mkbibnamefamily}[1]{%
  \ifitemannotation{highlight}{\textbf{#1}}{#1}}
\renewcommand*{\mkbibnamegiven}[1]{%
  \ifitemannotation{highlight}{\textbf{#1}}{#1}}
\renewcommand*{\mkbibnameprefix}[1]{%
  \ifitemannotation{highlight}{\textbf{#1}}{#1}}

% Hide the [1] [2] numeric labels — for a CV-style bibliography we just
% want hanging-indent entries with no labels.
\defbibheading{none}{}
\renewcommand{\bibsetup}{}
\setlength{\bibitemsep}{2pt plus 0.5pt minus 0.3pt}
\setlength{\bibhang}{1.2em}
\DeclareFieldFormat{labelnumberwidth}{}
\setlength{\biblabelsep}{0pt}
\defbibenvironment{bibliography}
  {\list{}{\setlength{\leftmargin}{\bibhang}%
           \setlength{\itemindent}{-\bibhang}%
           \setlength{\itemsep}{\bibitemsep}%
           \setlength{\parsep}{0pt}}}
  {\endlist}
  {\item}

% Drop the publist default "Talks" header noise; we only print bibliography.
\AtEveryBibitem{\clearfield{abstract}\clearfield{file}}

% --- Page styles ---------------------------------------------------------
% First page: empty (no header, no footer).
% Subsequent pages: small grey identifier top-right.
\fancypagestyle{cv}{
  \fancyhf{}
  \renewcommand{\headrulewidth}{0pt}
  \renewcommand{\footrulewidth}{0pt}
  \fancyfoot[L]{\sansfam\scriptsize\color{meta}Updated \today}%
  \fancyfoot[R]{\sansfam\scriptsize\color{meta}%
    Dennis Loevlie \,\textperiodcentered\, CV \,\textperiodcentered\, %
    page \thepage\ of \pageref*{LastPage}}
}
\fancypagestyle{cvfirst}{
  \fancyhf{}\renewcommand{\headrulewidth}{0pt}\renewcommand{\footrulewidth}{0pt}
  \fancyfoot[L]{\sansfam\scriptsize\color{meta}Updated \today}%
  \fancyfoot[R]{\sansfam\scriptsize\color{meta}%
    Dennis Loevlie \,\textperiodcentered\, CV \,\textperiodcentered\, %
    page \thepage\ of \pageref*{LastPage}}
}
\pagestyle{cv}

% --- Section / subsection ------------------------------------------------
% True small caps in serif body font, color accent, letter-spaced via
% microtype's textls. Thin accent rule under each section.
% Spaced uppercase Plex Sans (Tufte-style) — substitutes for the missing
% small caps in Plex's pdflatex variant. Letter-spacing brings it into
% the small-caps register. Tracking at 120 (~12%) per Butterick's
% "5-12% for all-caps" guidance — 160 was shouting.
\newcommand{\sectiontitle}[1]{\textls[120]{\MakeUppercase{#1}}}
\titleformat{\section}
  [block]
  {\sansfam\footnotesize\bfseries\color{accent}}
  {}{0em}{\sectiontitle}[\vspace{1pt}{\color{rule}\titlerule[0.4pt]}]
% Section spacing as baseline multiples for grid consistency.
\titlespacing*{\section}{0pt}{0.6\baselineskip plus 0.05\baselineskip minus 0.05\baselineskip}{0.2\baselineskip}

\titleformat{\subsection}
  {\footnotesize\itshape\color{meta}}
  {}{0em}{}
\titlespacing*{\subsection}{0pt}{0.3\baselineskip}{0pt}

% --- Lists ---------------------------------------------------------------
% Refined dot bullet (period-centered) — lighter than \textbullet,
% sits closer to baseline of serif body
\setlist[itemize]{itemsep=0.5pt, topsep=1pt, parsep=0pt,
                  leftmargin=1.0em,
                  label={\color{accent}\small\textperiodcentered}}

% --- Macros --------------------------------------------------------------
\newcommand{\me}{\textbf{Dennis Johan Loevlie}\xspace}
\newcommand{\meshort}{\textbf{Dennis Loevlie}\xspace}

% Em-dash with thin spaces for parenthetical breaks (not en-dash)
\newcommand{\bdash}{\,\textemdash\,}

% Date column: only place sans appears in body. Tabular lining numerals
% so date columns align cleanly across rows.
% Date column gets lining tabular figures (uniform width, cap-height) so
% dates align column-perfectly across rows. Body still uses oldstyle.
\newcommand{\datestamp}[1]{{\sansfam\footnotesize\color{meta}\figureversion{lf}#1}}

% Role line: bold serif label on left, sans tabular date on right.
% \filbreak treats each role+sub+bullets as atomic — page breaks may
% only happen between roles, never inside one. Combined with \nopagebreak
% after the role and sub lines so the heading never separates from its
% first bullet.
\newcommand{\role}[2]{%
  \filbreak%
  \par\vspace{1pt}%
  \noindent\textbf{#1}\hfill\datestamp{#2}\par\nopagebreak}

\newcommand{\sub}[1]{\noindent\textit{\footnotesize\color{meta} #1}\par\nopagebreak}

\newcommand{\monolink}[2]{\href{#1}{\monofam\footnotesize #2}}

% Venue treatment for publications: spaced caps in accent color rather
% than bold. Per Bringhurst, bold + italic on the same logical unit
% (already-italic-title + bold-venue) is "two emphases at once."
% Spaced caps reads as the venue tier without doubling the stress.
\newcommand{\venue}[1]{{\textls[80]{\MakeUppercase{#1}}}}

% Inline acronyms — Bringhurst rule 3.2.2: spaced small caps for
% acronyms in body text. Plex pdflatex ships no real small caps, so
% we use the same letter-spaced uppercase workaround as section
% headers, but at an even lighter tracking so the acronym integrates
% with body color rather than competing with section headers.
\newcommand{\acro}[1]{\textls[80]{#1}}

% Publications: hanging indent so author lists wrap under the title
\newenvironment{publist}{%
  \begin{list}{}{%
    \setlength{\leftmargin}{1.2em}%
    \setlength{\itemindent}{-1.2em}%
    \setlength{\itemsep}{2pt}%
    \setlength{\parsep}{0pt}%
    \setlength{\topsep}{1pt}%
  }%
}{\end{list}}

\begin{document}
\thispagestyle{cvfirst}

%---------------------------------------------------------------------------
% Header
%---------------------------------------------------------------------------
\begin{center}
{\sansfam\fontsize{22pt}{26pt}\selectfont\bfseries\color{accent} Dennis Johan Loevlie}\\[2pt]
% Research-interests one-liner. Olah-style positioning — lets advisors
% and recruiters triangulate the candidate in 2 seconds.
{\footnotesize\itshape\color{meta}Multimodal foundation models for tabular data \,\textperiodcentered\, probabilistic deep learning}\\[3pt]
{\sansfam\footnotesize\color{meta}%
+1\,724\,841\,8769 \,\textperiodcentered\,
\href{mailto:dennis.loevlie@tufts.edu}{dennis.loevlie@tufts.edu} \,\textperiodcentered\,
\href{https://www.dennisloevlie.com}{dennisloevlie.com} \,\textperiodcentered\,
\href{https://scholar.google.com/citations?user=oGkEIYkAAAAJ}{Google Scholar} \,\textperiodcentered\,
\href{https://github.com/loevlie}{github.com/loevlie}}
\end{center}
\vspace{-4pt}{\color{rule}\rule{\linewidth}{0.5pt}}

%---------------------------------------------------------------------------
% Education
%---------------------------------------------------------------------------
\section{Education}

\role{Centrum Wiskunde \& Informatica \& University of Amsterdam}{Incoming, Aug 2026 \textendash{}}
PhD in Computer Science, \textit{ELLIS PhD Program}. Advisors: Madelon Hulsebos (TRLab, CWI) and Jan-Willem van de Meent (AMLab, UvA). Focus: multimodal tabular foundation models.

\role{Tufts University, Medford, MA}{Sep 2024 \textendash{} May 2026 (expected)}
M.S.\ in Computer Science. GPA: 3.95\,/\,4.00. \href{https://www.loevliedl.com/static/Portfolio/PDF/SSR_TSRPT.pdf}{Transcript}.

\role{Carnegie Mellon University, Pittsburgh, PA}{Sep 2019 \textendash{} Dec 2020}
M.S.\ in Chemical Engineering. GPA: 3.91\,/\,4.00.

\role{West Virginia University, Morgantown, WV}{Sep 2016 \textendash{} Aug 2019}
B.S.\ in Chemical Engineering, \textit{cum laude}, with Honors.

%---------------------------------------------------------------------------
% Publications
%---------------------------------------------------------------------------
\section{Publications}

% \nocite{*} pulls every entry from publications.bib into the bibliography
% even though we don't explicitly \cite{} them in body text — the standard
% pattern for a CV's publications-as-bibliography setup.
\nocite{*}

\subsection{Machine Learning}
\printbibliography[heading=none, keyword=ml]

\subsection{Chemistry \& Computational Science}
\printbibliography[heading=none, keyword=chem]

%---------------------------------------------------------------------------
% Research Experience
%---------------------------------------------------------------------------
\section{Research Experience}

\role{Tufts University \bdash{} Hughes Lab}{Aug 2024 \textendash{} Present}
\sub{Graduate Researcher with Prof.\ Michael Hughes}
\begin{itemize}
  \item Co-led the \acro{ML4H} 2025 study showing a \emph{generalization gap} between correlated and non-correlated multiple instance learning (\acro{MIL}) architectures, using a synthetic benchmark with a Bayes-optimal upper bound.
  \item Identified and resolved a quadratic-scaling bottleneck in the lab's \acro{MIL} codebase, enabling experiments on substantially larger 3D image datasets.
  \item Designed attention- and transformer-based \acro{MIL} pipelines on 3D \acro{MRI}/\acro{CT} for early indicators of dementia and stroke; trained on Tufts \acro{HPC} with multi-\acro{GPU} distributed setups.
  \item Developing a regularization method to encourage interpretable attention scores, and methods for supervised learning under noisy labels.
\end{itemize}

\role{Tufts University \bdash{} Sinapov Lab}{Jan 2025 \textendash{} Present}
\sub{Graduate Researcher with Prof.\ Jivko Sinapov}
\begin{itemize}
  \item Trained Qwen2.5-Coder-7B with Group Relative Policy Optimization (\acro{GRPO}) and a custom 3-part reward (structure, aesthetics, semantic alignment) to generate \acro{SVG}s from text; achieved an \textbf{18\% improvement} on a benchmark covering aesthetics, alignment, and code validity.
  \item Investigating why vision-language models (Qwen3-\acro{VL}) and open-vocabulary segmentation models (\acro{SAM}\,3) fail at object counting in dense scenes; designing tool-use and trajectory-based methods drawing on cognitive-science accounts of human counting.
\end{itemize}

\role{University of Pittsburgh \bdash{} CANELa Lab}{Jun 2021 \textendash{} Jan 2023}
\sub{Graduate Researcher with Prof.\ Giannis Mpourmpakis}
\begin{itemize}
  \item Proposed a weight-initialization method that delivered a \textbf{71\% \acro{RMSE} reduction} for predicting nanoparticle stability across compositions and atomic orderings (\emph{Loevlie et al., Acc.\ Chem.\ Res.\ 2023}).
  \item Designed and tuned \acro{ML} architectures for material-property prediction; co-author on \emph{Salem et al., JPCC 2023}.
  \item Wrote the \acro{ML} methods and background section for the \emph{Nagarajan et al.}\ review (Curr.\ Opin.\ Chem.\ Eng., 2022).
\end{itemize}

\role{Carnegie Mellon University \bdash{} Kitchin Group}{Dec 2019 \textendash{} Dec 2020}
\sub{Graduate Researcher with Prof.\ John Kitchin}
\begin{itemize}
  \item Trained a \acro{CNN} classifier to extract experimental information from microscopy images; rebuilt Mathematica analysis tooling as fast, interactive Python.
  \item Released \href{https://tinyurl.com/yzpperdc}{\monofam nb\_search}, an open-source Python package for searching and indexing across Jupyter notebooks.
\end{itemize}

%---------------------------------------------------------------------------
% Awards
%---------------------------------------------------------------------------
\section{Awards \& Honors}

\begin{list}{}{%
  \setlength{\leftmargin}{3.4em}%
  \setlength{\labelwidth}{2.6em}%
  \setlength{\labelsep}{0.5em}%
  \setlength{\itemsep}{1pt}%
  \setlength{\parsep}{0pt}%
  \setlength{\topsep}{1pt}}
\item[\datestamp{2026}] \textbf{1st place} (10 teams), \textit{BrainStorm Neural Decoder Challenge}\bdash{}real-time auditory decoding from 1024-channel \acro{ECoG}; 95\% accuracy, sub-ms inference, edge-deployable model.
\item[\datestamp{2024}] \textbf{Community Grant}, \textit{Hugging Face}\bdash{}for \href{https://depth-anything.github.io/}{Depth Anything} demonstration work on video.
\item[\datestamp{2019}] \textbf{1st place} (Advanced category), \textit{AVEVA National Simulation Competition}.
\item[\datestamp{2018}] \textbf{2nd place}, \textit{AIChE National Poster Competition}, Computing \& Process Control Division.
\end{list}

%---------------------------------------------------------------------------
% Industry Experience
%---------------------------------------------------------------------------
\section{Industry Experience}

\role{Massachusetts General Hospital \& Harvard Medical School}{2025 \textendash{} Present}
\sub{Machine Learning Researcher \bdash{} computational pathology}
\begin{itemize}
  \item \textbf{Systems:} building segmentation models for colon-cancer pathology slides at gigapixel resolution; pipeline integrates with clinical workflows for downstream evaluation.
  \item \textbf{Research:} dataset-construction methodology and collaboration with pathologists to ground model outputs in expert-validated labels.
\end{itemize}

\role{KEF Robotics, Pittsburgh, PA}{Jan 2023 \textendash{} Aug 2024}
\sub{Senior Computer Vision \& \acro{ML} Engineer (2024) \,$\cdot$\, Computer Vision Engineer (2023)}
\begin{itemize}
  \item \textbf{Systems:} sole \acro{ML} engineer on a one-year, \$500K project; led a team of five shipping on-device object detection, monocular depth, and 3D map generation; deployed on flight hardware via TensorRT; \textbf{45\% inference-speed gain} at $<\!1\%$ accuracy cost. Two in-person demos, selected for follow-on funding.
  \item \textbf{Research:} fine-tuned Mask2Former for power-line segmentation across \acro{RGB} and \acro{IR} modalities via transfer learning and cross-modal label transfer.
\end{itemize}

\role{ContentsPal}{May 2025 \textendash{} Aug 2025}
\sub{Multimodal \acro{AI} Intern (insurance-tech startup, \acro{MIT}-led)}
\begin{itemize}
  \item \textbf{Systems:} integrated and tested learning-based duplicate detection and open-vocabulary instance segmentation in React / React Native.
  \item \textbf{Research:} prototyped retrieval and matching components on multimodal product data.
\end{itemize}

\role{AiThElite}{Dec 2020 \textendash{} May 2021}
\sub{Lead Data Scientist (collegiate-athletics AI startup)}
\begin{itemize}
  \item Built end-to-end data pipeline (Selenium scrapers \textrightarrow{} NumPy/Pandas feature engineering \textrightarrow{} scikit-learn models) and shipped a Django/AWS-hosted prediction product.
\end{itemize}

%---------------------------------------------------------------------------
% Open Source
%---------------------------------------------------------------------------
\section{Open Source}

\begin{itemize}
  \item \href{https://github.com/huggingface/transformers}{\monofam huggingface/transformers} \,\textemdash\, contributor; merged PRs on model implementations and bug fixes
  \item \href{https://github.com/Franblueee/torchmil}{\monofam Franblueee/torchmil} \,\textemdash\, core contributor; multiple merged PRs across the multiple-instance-learning library
  \item \href{https://github.com/loevlie/neuropt}{\monofam loevlie/neuropt} \,\textemdash\, author; \acro{LLM}-guided neural architecture search and hyperparameter tuning
  \item \href{https://huggingface.co/spaces/JohanDL/Depth-Anything-Video}{\monofam JohanDL/Depth-Anything-Video} \,\textemdash\, author; Hugging Face Community Grant recipient
\end{itemize}

%---------------------------------------------------------------------------
% Talks
%---------------------------------------------------------------------------
\section{Selected Talks}

\begin{publist}
\item \href{https://drive.google.com/file/d/1Z6lkTsqYSIzOqgM53qONW2riQhVM5R0W/view?usp=sharing}{\textit{Synthetic Data Reveals Generalization Gaps in Correlated Multiple Instance Learning}}. Poster, ML4H 2025 Symposium, San Diego, CA. \textit{Dec 2025.}

\item \href{https://xchangepgh2023.sched.com/event/1LdC7/an-interactive-exploration-of-computer-vision-for-unmanned-aerial-vehicles-uavs}{\textit{Computer Vision for Unmanned Aerial Vehicles}}. Invited talk, XChangeIdeas Pittsburgh. \textit{2023.}
\end{publist}

%---------------------------------------------------------------------------
% Service & Outreach
%---------------------------------------------------------------------------
\section{Service \& Outreach}

Industry volunteer, youth robotics team building accessibility tools for blind soccer players \hfill \datestamp{2024}

%---------------------------------------------------------------------------
% Skills
%---------------------------------------------------------------------------
\section{Technical Skills}

\noindent\textbf{Languages} \quad Python $\cdot$ C++ $\cdot$ JavaScript $\cdot$ MATLAB \\
\textbf{\acro{ML} / AI} \quad PyTorch $\cdot$ JAX $\cdot$ Hugging Face (Transformers, Datasets, \acro{TRL}, \acro{PEFT}) $\cdot$ \acro{GRPO} / \acro{DPO} $\cdot$ multimodal learning $\cdot$ \acro{RL} post-training $\cdot$ multiple instance learning $\cdot$ computer vision \\
\textbf{Engineering} \quad \acro{FSDP} $\cdot$ DeepSpeed $\cdot$ vLLM $\cdot$ FlashAttention $\cdot$ Triton $\cdot$ \acro{CUDA} $\cdot$ TensorRT $\cdot$ \acro{SLURM} / \acro{HPC} $\cdot$ multi-\acro{GPU} distributed training $\cdot$ Docker $\cdot$ \acro{AWS}

\end{document}
